% AUTO-GENERATED by scripts/exp_real_input_40_event_time_rescaling_error_budget.py
\begin{table}[H]
\centering
\scriptsize
\setlength{\tabcolsep}{6pt}
\renewcommand{\arraystretch}{1.15}
\caption{事件换时的期望有效指数与偏移：定义 $\tau^{(E)}(k):=\log\mathbb{E}(\rho^{\tau(k)})/\log\rho$，并记 $\Delta^{(E)}(k):=\tau^{(E)}(k)-k/\mu_e$（真实输入 40 状态核；$t=0$ Parry 平衡态）。设 $\mu_e=0.367376207875$，$\rho=\rho_{\mathrm{corr}}(0)=0.618033988750$，$\tau_{\mathrm{mix}}=19.747340579192$，并用 $n_{\max}=260$ 的有限递推与尾界得到 $\mathbb{E}(\rho^{\tau(k)})$ 的严格区间（见表 \ref{tab:real-input-40-event-time-rho-tau-expectation-vs-tau-mix}）。由 $\log$ 单调性与 $\log\rho<0$，可把该区间推出 $\tau^{(E)}(k)$ 与 $\Delta^{(E)}(k)$ 的严格区间。}
\label{tab:real-input-40-event-time-tauE-delta-vs-tau-mix}
\begin{tabular}{r r r r r r}
\toprule
$k$ & $k/\mu_e$ & $\tau^{(E)}(k)_{\mathrm{lo}}$ & $\tau^{(E)}(k)_{\mathrm{hi}}$ & $\Delta^{(E)}(k)_{\mathrm{lo}}$ & $\Delta^{(E)}(k)_{\mathrm{hi}}$\\
\midrule
6 & 16.332032046126 & 14.353900876823 & 14.353900876823 & -1.978131169304 & -1.978131169304\\
7 & 19.054037387147 & 16.894564896219 & 16.894564896219 & -2.159472490928 & -2.159472490928\\
8 & 21.776042728168 & 19.435250893884 & 19.435250893884 & -2.340791834284 & -2.340791834284\\
9 & 24.498048069189 & 21.975944888338 & 21.975944888338 & -2.522103180851 & -2.522103180851\\
10 & 27.220053410210 & 24.516641780697 & 24.516641780697 & -2.703411629513 & -2.703411629513\\
\bottomrule
\end{tabular}
\end{table}
